\chapter{Конструкторская часть}

\section{Требования к программному обеспечению}

К разрабатываемой программе предъявлен ряд требований:

\textbf{Входные данные:} Массив целых чисел, искомое целое число.

\textbf{Выходные данные:} Индекс искомого числа в массиве.

\begin{itemize}
	\item Индексация элементов в массиве начинается с 0;
	\item В случае, если искомого элемента в массиве нет, вместо индекса должно возвращаться значение -1.
\end{itemize}

\section{Вариации бинарного поиска}

Алгоритм бинарного поиска можно реализовать двумя способами: итеративно и рекурсивно.

В итеративной версии используются две переменные для определения левой и правой границ. В цикле вычисляется средний элемент, и границы обновляются в зависимости от его значения.

В рекурсивной версии определяется средний элемент всего массива, и в зависимости от его значения рекурсивно вызывается поиск либо в правой, либо в левой части массива.

В данной работе будет рассматриваться итеративная реализация алгоритма.

\clearpage

\section{Разработка алгоритмов}

Схема алгоритма поиска в массиве полным перебором представлена на рисунке \ref{fig:src}.

\begin{figure}[H]
	\centering
	\includegraphics[width=0.95\textwidth]{src}
	\caption{Схема алгоритма поиска в массиве полным перебором}
	\label{fig:src}
\end{figure}

\clearpage

Схема алгоритма бинарного поиска представлена на рисунке \ref{fig:bin}.

\begin{figure}[H]
	\centering
	\includegraphics[width=0.95\textwidth]{bin}
	\caption{Схема алгоритма бинарного поиска}
	\label{fig:bin}
\end{figure}

\section{Вывод}

В результате конструкторской части были определены требования к ПО, спроектированы алгоритмы поиска в массиве полным перебором и бинарного поиска.
